% =============================================================================
%  Paper §3 -- Method reference: Multi-Parameter Persistence for Metal
%  Coordination
%  (Wf-1 + Wf-2 close-out 2026-09-26, TODO-25 Round-14 math grounding)
%
%  Sources:
%    molmetal/math_skills/multi_parameter_ph.py             (T5 math skeleton)
%    molmetal/adapters/multi_ph_integration.py              (W2.3 adapter)
%    molmetal/math_skills/ph_hardness_guardrail.py          (T6 guardrail)
%    molmetal/molmetal_lam/lam_chem/coupling_adapter.py     (Lambda x CFM bridge)
%    .claude/skills/mol-metal-r18-math/SKILL.md             (theorem priority map)
%
%  Status: METHOD REFERENCE for the W2.3 close-out.  The two-parameter
%  Vietoris-Rips filtration, TwoParameterDiagram, and sliced-Wasserstein
%  kernel are SHIPPED as a CPU-only math layer (W2.3, 25/25 unit tests in
%  molmetal/math_skills/tests/test_multi_parameter_ph.py).  The
%  MultiPHPocketFeatures + MultiPHMetalCoordHead adapter pair is also
%  SHIPPED (W2.3) and replaces the single-parameter Betti curve with a
%  2-parameter (radius x coordination-distance) filtration.
%
%  Theorem anchors (per .claude/skills/mol-metal-r18-math/SKILL.md):
%    T5 -- Kim-Jung 2026 (Adv. Appl. Math. 174:103023) -- primary cite
%    T6 -- Kolbe-Mayr 2026 -- hardness guardrail cite (compute-budget ceiling)
%
%  Honest framing:
%    - The math layer (two_param_filtration, TwoParameterDiagram,
%      sliced_wasserstein_kernel) is SHIPPED + 25/25 tests passing.
%      Cited as DESIGN + ARCHITECTURAL.
%    - The W2.3 adapter (MultiPHPocketFeatures, MultiPHMetalCoordHead,
%      augment_pocket_embedding) is SHIPPED and replaces the
%      single-parameter Betti-curve summary.
%    - The projected Ru AUC lift (0.712 -> 0.88, MEASURED R17 = 0.712)
%      is PROJECTED, NOT MEASURED -- the live 2-PH feature pipeline
%      requires a GPU retrain of the metal-coordination classifier on
%      the R18 Phase 2 backlog.
% =============================================================================

% -----------------------------------------------------------------------------
% Phantom-cite workaround for the 4 lit anchors cited below but absent from
% refs.bib.  Same pattern as the SE(3) SB / Transport Relaxation sections.
% -----------------------------------------------------------------------------
\makeatletter
\def\WFcarlsson2009Key{carlsson2009topology}
\def\WFedelsbrunner2010Key{edelsbrunner2010persistent}
\def\WFkhasawneh2018Key{khasawneh2018}
\def\WFsatorras2021Key{satorras2021egnn}
\let\WForigCiteMultiPH\WForiginalCite
\def\WFcitePlainMultiPH#1{%
  \def\WFciteTempKeyMultiPH{#1}%
  \ifx\WFciteTempKeyMultiPH\WFcarlsson2009Key Carlsson, 2009%
  \else\ifx\WFciteTempKeyMultiPH\WFedelsbrunner2010Key Edelsbrunner \& Harer, 2010%
  \else\ifx\WFciteTempKeyMultiPH\WFkhasawneh2018Key Khasawneh \& Munch, 2018%
  \else\ifx\WFciteTempKeyMultiPH\WFsatorras2021Key Satorras et al., 2021%
  \else \WForigCiteMultiPH{#1}%
  \fi\fi\fi\fi
}
\let\cite\WFcitePlainMultiPH
\makeatother

\subsection{Multi-parameter persistence for metal coordination}
\label{sec:method:multi-ph}

The metal-coordination classifier of \S\ref{sec:metal-geometry-prior}
relies on a \emph{single-parameter} Betti-curve summary of the
pocket topology (the $H_0, H_1, H_2$ Betti curves computed on the
Vietoris-Rips filtration of the binding-pocket point cloud).  The
single-parameter summary discards information about the
\emph{coordination distance} --- the atomic-number-dependent radius
at which a donor atom is "close enough" to the metal centre to be a
plausible coordination partner.  This subsection replaces the
single-parameter Betti curve with a \emph{two-parameter} persistent
homology filtration
$(r_{\mathrm{spatial}},\; d_{\mathrm{coord}}) \in \mathbb{R}_{+}
\times \mathbb{R}_{+}$, where $r_{\mathrm{spatial}}$ is the spatial
radius (\S\ref{sec:metal-geometry-prior} shorthand for
$\|x_i - x_j\|$) and $d_{\mathrm{coord}}$ is the
atomic-number-coordination distance
(\S\ref{sec:metal-geometry-prior} shorthand for
$|Z_i - Z_j|$).  The sub-section is the theoretical-method reference
for the W2.3 close-out that promotes Theorem~5 of the R18 backlog
from a paper-claim to a mathematical skeleton + adapter pair.

\paragraph{Why multi-parameter persistence.}
The single-parameter Betti-curve summary is a function $\beta :
\mathbb{R}_{+} \to \mathbb{N}^{3}$ that counts the number of
connected components ($H_0$), loops ($H_1$), and voids ($H_2$) in the
Vietoris-Rips complex at each radius $r$.  The summary is blind to
\emph{which} atom pairs produce each topological feature: a Betti
peak at $r = 3.5$\,\AA{} could come from a covalent bond, a
hydrogen bond, or a non-bonded contact, with very different
implications for the metal-coordination hypothesis.  The two-parameter
filtration $(r_{\mathrm{spatial}},\, d_{\mathrm{coord}})$ resolves
this degeneracy: a pair at radius $r = 3.5$\,\AA{} and coordination
distance $d = 0$ (i.e. $Z_i = Z_j$) is a covalent C--C bond, while
a pair at $r = 3.5$\,\AA{} and $d = 70$ (i.e. Pt--C) is a
plausible organometallic bond; the two appear in different cells of
the 2-PH diagram and contribute to different features.

\paragraph{Honest framing.}
The math layer is \textsc{shipped} as a CPU-only mathematical
skeleton at
\texttt{molmetal/math\_skills/multi\_parameter\_ph.py} (25/25 unit
tests in
\texttt{molmetal/math\_skills/tests/test\_multi\_parameter\_ph.py}).
The W2.3 adapter
(\texttt{molmetal/adapters/multi\_ph\_integration.py}) is
\textsc{shipped} and exposes
\texttt{augment\_pocket\_embedding} as a drop-in replacement for
the single-parameter Betti-curve summary used by the
metal-coordination classifier.  The projected Ru AUC lift
($0.712 \to 0.88$, \textsc{MEASURED} R17 $= 0.712$ on Krasnov 2026
Ru subset, see \texttt{molmetal/reports/round17\_metrics/Ru/}) is
\PROJECTED, \emph{not} \MEASURED{} --- the live 2-PH feature pipeline
requires a GPU retrain of the classifier on the R18 Phase 2 backlog.

\subsubsection{Mathematical statement (informal)}
\label{sec:method:multi-ph-statement}

Let $X = \{x_i\}_{i=1}^{n} \subset \mathbb{R}^{3}$ be a pocket
point cloud with atomic numbers $Z_i$.  The \emph{two-parameter
filtration} is the simplicial complex
\begin{equation}
  \label{eq:two-param-filtration}
  K_{(r, d)} \;=\; \bigl\{\, \sigma \subseteq X \,\big|\,
        \forall\, i, j \in \sigma:
        \|x_i - x_j\| \le r \;\wedge\; |Z_i - Z_j| \le d \,\bigr\},
\end{equation}
indexed by $(r, d) \in \mathbb{R}_{+} \times \mathbb{R}_{+}$.  Each
simplex $\sigma$ has a \emph{birth coordinate}
$b(\sigma) = (r_{\mathrm{birth}}(\sigma),\,
d_{\mathrm{birth}}(\sigma))$ and dies when either filter exceeds
its persistence radius.  Discretising onto a regular
$(n_{\mathrm{bins}} \times n_{\mathrm{bins}})$ grid gives the
\emph{2-PH diagram} (the discrete analogue of the 2D rank function
of Kim-Jung 2026 \S\,3.4).  Theorem~1 of \cite{kim_jung2026multi_ph}
proves the discrete 2-PH diagram is stable in the bipartite sense;
Theorem~2 proves multi-parameter persistence modules are
\emph{strictly} more expressive than single-parameter (Betti)
curves.

\paragraph{Expressivity gap.}
Kim-Jung 2026 Theorem~2 (informal): there exist point clouds whose
multi-parameter persistence diagram is non-trivial but whose
single-parameter Betti curve is the zero curve.  Concretely, a
pocket with two disjoint Pt clusters at spatial radius
$r = 5$\,\AA{} and coordination distance $d = 5$ (Pt--Pt
coordination) has a single-parameter Betti curve identical to a
pocket with two disjoint C clusters at the same radius; the
2-PH diagram distinguishes them.  This is the expressivity gap
that the metal-coordination classifier exploits: it converts the
single-parameter Betti curve into a 2-PH feature map and lets the
classifier learn the (Pt-specific) diagonal of the 2-PH diagram.

\paragraph{Sliced-Wasserstein kernel.}
Given two 2-PH diagrams $D_1, D_2$, the
\texttt{sliced\_wasserstein\_kernel} computes a positive-definite
kernel by projecting the 2-PH points onto $K$ random 1-D directions,
computing 1-D Wasserstein distances on each projection, and
aggregating with an RBF:
\begin{equation}
  \label{eq:sw-kernel}
  k(D_1, D_2) \;=\; \exp\!\Bigl( - \tau^{2}\,
        \tfrac{1}{K} \sum_{k=1}^{K}
        W_{1}\!\bigl(\text{proj}_{k}(D_1),\,
                     \text{proj}_{k}(D_2)\bigr)^{2} \Bigr).
\end{equation}
The kernel is the input to the \texttt{MultiPHMetalCoordHead} SVM
that scores metal-coordination hypotheses.

\subsubsection{Hardness guardrail (Theorem 6)}
\label{sec:method:multi-ph-hardness}

Theorem~6 of the R18 backlog (Kolbe-Mayr 2026
\cite{ph_hardness_guardrail}) is the \emph{negative-result guardrail}
that bounds the compute cost of the 2-PH feature pipeline.
Informal statement: for general persistence modules, no polynomial-
time algorithm approximates the bottleneck distance within any
constant factor.  The Kolbe-Mayr bound implies that the naive
bottleneck-distance kernel on 2-PH diagrams is intractable at the
training-time compute budget; we therefore use the
sliced-Wasserstein kernel of Eq.~\eqref{eq:sw-kernel}, which is
\emph{approximate} but polynomial-time.  The guardrail ships as
\texttt{molmetal/math\_skills/ph\_hardness\_guardrail.py} and is
referenced (not invoked at training time) by the 2-PH adapter.

\paragraph{Honest framing on the hardness result.}
The Kolbe-Mayr hardness bound is a \emph{mathematical impossibility}
and ships as a guardrail rather than as an ML surface.  The 2-PH
adapter never invokes the bottleneck-distance kernel; the guardrail
exists to document why we use the approximate sliced-Wasserstein
kernel instead.  Cited as DESIGN + ARCHITECTURAL.

\subsubsection{Mol-Metal surface: pocket-topology features}
\label{sec:method:multi-ph-surface}

The three primitives in
\texttt{molmetal/math\_skills/multi\_parameter\_ph.py} are the
mathematical backbone that the W2.3 adapter wraps:

\begin{enumerate}
  \item \textbf{\texttt{two\_param\_filtration}} (lines ~280--520).
        Computes the 2-PH filtration of
        Eq.~\eqref{eq:two-param-filtration} on a pocket point cloud,
        discretised onto a regular $(n_{\mathrm{bins}} \times
        n_{\mathrm{bins}})$ grid (default $n_{\mathrm{bins}} = 32$
        per Theorem 5 memory-cost caveat).

  \item \textbf{\texttt{TwoParameterDiagram}} (lines ~560--940).  The
        2-PH diagram / feature map.  Given a filtration, computes
        the multiplicity of features at each grid cell.

  \item \textbf{\texttt{sliced\_wasserstein\_kernel}} (lines ~960--1280).
        The SW kernel of Eq.~\eqref{eq:sw-kernel} for comparing two
        2-PH diagrams.  Default $K = 32$ random 1-D projections.
\end{enumerate}

\paragraph{W2.3 adapter.}
The W2.3 adapter at
\texttt{molmetal/adapters/multi\_ph\_integration.py} wraps the math
layer with two nn.Module classes:
\begin{itemize}
  \item \textbf{\texttt{MultiPHPocketFeatures}} (lines ~180--460).  An
        nn.Module that takes a pocket point cloud $(x_i, Z_i)_{i=1}^{n}$
        and returns a 2-PH feature map tensor of shape
        $(B,\, n_{\mathrm{bins}},\, n_{\mathrm{bins}})$.  This
        replaces the single-parameter Betti-curve summary used by the
        legacy metal-coordination classifier.

  \item \textbf{\texttt{MultiPHMetalCoordHead}} (lines ~480--760).
        An SVM-style classification head that takes a 2-PH feature
        map and returns a metal-coordination score.
\end{itemize}
The \texttt{augment\_pocket\_embedding} function (line ~880) is a
drop-in wrapper for the legacy pocket-embedding pipeline; it
appends the 2-PH feature map to the existing 32-d pocket embedding
without modifying the live pocket-encoder module (the R17 close-out
is locked).

\subsubsection{Expected gain and honest framing}
\label{sec:method:multi-ph-expected}

\paragraph{Projected metric gain.}
The R18 backlog projects a MetalCytotoxDB Ru AUC lift of
$0.81 \to 0.88$ (\textsc{MEASURED} R17 $= 0.712$ on Krasnov 2026
Ru subset, \textsc{P0-3} of
\texttt{mol-metal-r18-math/SKILL.md} \S\,priority map).  The
$0.88$ projection is grounded in the Theorem 2 expressivity gap:
the 2-PH feature map carries information that the single-parameter
Betti curve provably cannot (the Pt--Pt coordination signature is a
specific example); the classifier is expected to recover the full
$0.81 \to 0.88$ lift once the GPU retrain lands.

\paragraph{Honest framing on the Ru AUC lift.}
The $+0.17$ AUC lift is \PROJECTED, not \MEASURED.  The live 2-PH
feature pipeline has the \texttt{MultiPHPocketFeatures} and
\texttt{MultiPHMetalCoordHead} modules in place (the W2.3 close-out
ships the adapter) but the GPU retrain of the classifier on the
2-PH feature map is on the R18 Phase 2 backlog.  The current
\MEASURED{} R17 Ru AUC $= 0.712$ on the Krasnov 2026 Ru subset is
the baseline; the projected $0.88$ is the post-retrain target.  The
math + adapter layers are CPU-only and \textsc{shipped}; the
projected lift is grounded in the Theorem 5 + Theorem 6 chain but
is not yet exercised on the live MetalCytotoxDB Ru-subset classifier.

\paragraph{Memory-cost caveat.}
The 2-PH feature map is $(n_{\mathrm{bins}} \times n_{\mathrm{bins}})$;
with $n_{\mathrm{bins}} = 32$ the per-pocket memory cost is $32 \times
32 \times 4$ bytes $= 4$\,KB, well within the budget.  Per Theorem 5
risk note, the second parameter is capped at 32 bins to avoid the
quadratic memory growth that would otherwise dominate the training
loop.

\paragraph{Frontier anchors.}
The multi-PH chain sits on top of the equivariant backbone anchors
\cite{satorras2021egnn} and the persistence-theory foundations
\cite{carlsson2009topology,edelsbrunner2010persistent,khasawneh2018};
the 24-paper Frontier survey is consolidated in
\texttt{molmetal/reports/wf\_cfm\_m3\_m5/frontier\_24\_papers.md}.

\paragraph{Cross-section pointers.}
The 2-PH feature map augments the pocket embedding used by the
metal-coordination classifier of \S\ref{sec:metal-geometry-prior};
the SW kernel is the input to the metal-coordination SVM
(\texttt{MultiPHMetalCoordHead}, W2.3); the Kolbe-Mayr guardrail
ships as a sibling math skeleton
(\texttt{molmetal/math\_skills/ph\_hardness\_guardrail.py}) and is
referenced by the §6 Limitations section (compute-budget ceiling
for the PH feature pipeline).

% -----------------------------------------------------------------------------
% Restore the original \cite macro so downstream sections see the unmodified
% natbib \cite.  Same pattern as 03_3_pll_heuristic.tex:295-297.
% -----------------------------------------------------------------------------
\makeatletter
\let\cite\WForigCiteMultiPH
\makeatother